\vspace{-3mm}
\section{Related Works}
\vspace{-3mm}

\paragraph{Nonparametric IV problem.} 
% \masa{Can be much more concise by focusing on comparing to our work. Generally, we should't move the related work to Appendix because it is important to show our novelty. }
Nonparametric IV estimation has been extensively explored in past decades. Such estimation is tough to solve even when both the linear operator $\Tcal$ and the response $r_0$ are known, known as ill-posedness. The ill-posedness often refers to the presence of one or more
of the following characteristics: (1) the absence of solutions, (2) the existence of multiple solutions, and (3) the discontinuity of the inverse of operator $\Tcal$. Many traditional nonparametric estimators have been proposed to address these challenges, such as series-based estimators \citep{florens2011identification,ai2003efficient,chen2021robust,chen2012estimation,darolles2011nonparametric} and kernel-based estimators \citep{hall2005nonparametric,horowitz2007asymptotic,singh2019kernel}. However, these methods cannot directly accommodate modern machine-learning techniques like neural networks.

Recently, there has been growing interest in the application of general function approximation techniques, such
as deep neural networks and random forests, to IV problems in a unified manner. Among those methods, \citet{bennett2020variational,dikkala2020minimax,lewis2018adversarial, liao2020provably,zhang2023instrumental} reformulate the conditional moment constraint into a minimax optimization and use its solution as the estimator. Notably, \citet{liao2020provably,bennett2023source,bennett2023minimax} establish $L_2$ convergence by linking minimax optimization with Tikhonov regularization under the assumption of the source condition. Moreover, \citet{liao2020batch} assumes uniqueness of solution $h_0$. \cite{dikkala2020minimax,lewis2018adversarial} provide a guarantee for the projected MSE without further assumptions. However, they could not guarantee the convergence rate in strong $L_2$ metric when multiple solutions to conditional moment constraint exist. Furthermore, these methods require a computation oracle for minimax optimization, which further makes model selection challenging. 
%%%further hinders performing model selection when multiple function classes are involved. %%%%Indeed, it is still unknown how to perform model selection given this minimax regression procedure. 
In contrast, our method does not require computational oracles and enables model selection with statistical guarantees. 

Several existing works eschew the need for minimax optimization oracles \citep{hartford2017deep, xu2021deep}. However, all these works do not provide finite sample guarantee or model selection. For example, as the most related work, DeepIV \citep{hartford2017deep} introduces a similar loss function to us. However, it 
lacks an explicit regularization term, which results in the lack of theoretical guarantee and the lack of guarantee for model selection. As another work, \cite{xu2021deep} extends the two-stage kernel algorithm in \cite{singh2019kernel} to deep neural networks, but their algorithm is essentially a bilevel optimization problem, which is hard to solve in general \citep{hong2023two,khanduri2021near,guo2021randomized}. 

\paragraph{Model selection.} 

Model selection has been well studied in the regression and supervised machine
learning literature \citep{bartlett2002model,gold2003model,mcallester2003pac}. The objective can be described more concretely as follows: given $M$ candidate models, $\{f_1,\dots,f_M\}$, each having some statistical complexity $\delta_j$ and some approximation error $\epsilon_j$ (with respect to some un-known true model $f_0$) we wish to find an aggregated model $\hat{f}$ whose mean squared error is closed to the optimal trade-off between statistical complexity and approximation error among all models, i.e.:
$\|\hat{f} - f_0\| \lesssim \min_{j=1}^M \delta_j + \epsilon_j.$
The statistical complexity of a function space can be accurately characterized, albeit the approximation error is un-attainable as it relates to the unknown true model. A guarantee of the form above implies that using the observed data we can compete (up to constants) with an oracle that knows the approximation errors and chooses the best model space. We leave the detailed summary of existing works in Appendix \ref{sec:related-works}. Despite the abundance of methodologies for IV regression problems, few studies have investigated model misspecification and provided model selection procedures to select the best model class. As a few exceptional works, while \cite{xu2021deep} and \cite{AI20075} considered the misspecified regime, but they did not discuss model selection approaches. A typical approach to model selection is out-of-sample validation: estimate different models on half the data and select the estimated model that achieves the smallest empirical risk on the second half (or the best convex ensemble of models that achieves the smallest out-of-sample risk). One problem that arises for model selection in this IV regression setup is to transform the excess risk guarantees, which will be in terms of the weak metric, i.e. $\|\Tcal (\cdot)\|_2$, into the desired bound in the $L_2$ error. In this work, we show that by leveraging the Tikhonov regularization, we can achieve an MSE bound that achieves the same order as the oracle function class. 